% Siklus SDLC metode prototyping (TikZ).
% Di-input dari BAB III §3.1.2.1.

\begin{figure}[H]
\centering
\resizebox{0.72\textwidth}{!}{%
\begin{tikzpicture}[
  phase/.style={draw=black, fill=white, rectangle, rounded corners=4pt,
    minimum width=3.2cm, minimum height=1.0cm, align=center, font=\footnotesize\bfseries},
  center/.style={draw=none, fill=none, align=center, font=\small\itshape},
  ar/.style={-{Stealth}, very thick},
]

%--- Empat fase membentuk lingkaran ---
%  0°=kanan, 90°=atas  →  posisi fase:
%  Atas: Analisis Kebutuhan
%  Kanan: Perancangan
%  Bawah: Implementasi
%  Kiri: Evaluasi

\node[phase] (req)   at ( 0,  3.4) {Analisis\\Kebutuhan};
\node[phase] (design) at ( 3.4, 0) {Perancangan\\Purwarupa};
\node[phase] (impl)  at ( 0, -3.4) {Implementasi};
\node[phase] (eval)  at (-3.4, 0) {Evaluasi \&\\Pengujian};

\node[center] at (0, 0) {Siklus\\\textit{Prototyping}};

%--- Panah siklus (searah jarum jam) ---
\draw[ar] (req.east)   to[bend left=25]  (design.north);
\draw[ar] (design.south) to[bend left=25] (impl.east);
\draw[ar] (impl.west)  to[bend left=25]  (eval.south);
\draw[ar] (eval.north) to[bend left=25]  (req.west);

%--- Label panah ---
\node[font=\scriptsize, align=center] at ( 2.8,  2.4) {rancang\\\textit{prototype}};
\node[font=\scriptsize, align=center] at ( 2.8, -2.4) {bangun\\\textit{prototype}};
\node[font=\scriptsize, align=center] at (-2.8, -2.4) {uji \&\\evaluasi};
\node[font=\scriptsize, align=center] at (-2.8,  2.4) {perbaiki\\kebutuhan};

\end{tikzpicture}%
}
\caption{Siklus SDLC dengan metode \textit{prototyping}: empat fase berulang hingga purwarupa memenuhi kebutuhan fungsional dan kinerja}
\label{fig:sdlc-proto}
\end{figure}
